\scanpage{101}

\uline{N.B.} Il y aurait lieu d'examiner aussi la situation où
on a un faisceau loc. constant
$G$ \uline{sur $V$} \uline{seulement} [pas néc.
prolongeable à $X$], \deleted{\unsure{Du fait alors}} et qu'on
se propose d'étudier \inserted{les} $\mathbb R^i h_*(G)$.

On a une suite \deleted{\unsure{exacte}} spectrale

\[ \mathbb R^{\unsure{*|n|i}}\unsure{h_*|h}(G)\Longleftarrow \mathbb R^p f_*(\mathbb R^q i_*(G)) \]

où $i:V\unsure{-|\to}X$ est l'immersion.

Tout est bien si les $\mathbb R^p f_*(\mathbb R^q i_*(G))$
sont loc. constants. Le Pb. ne
semble pas facile, \inserted{même si $Y$ lisse$/S$,} à cause des
phénomènes de ``wild ramification''
qui peuvent se produire [si $\operatorname{codim}(Y,X)$
$\geq 2$, il n'y a pas d'ennui, car $G$
se \unsure{prolonge} alors en un faisceau
loc. constant \ldots\thinspace.]. \deleted{Peut-être \unsure{est}}

% archive sheet number: 100
